\section{Introduction}
\label{sec:intro}

A language model does one thing: given a sequence of tokens, it produces a distribution over the next token. It cannot see a file it was not shown, remember a conversation it was not sent again, or carry out a command it wrote. Language models are nonetheless described as agents that read codebases, run tests, browse the web and fix bugs. The difference between what a model is and what an agent does is software around the model. That software is the subject of this paper.

The software has many names. Evaluations call it a \emph{scaffold} \citep{anthropic2025swebench,metr2025update}; libraries call themselves \emph{frameworks}, \emph{SDKs} or an \emph{orchestration framework and runtime} \citep{langgraph}; and \emph{harness} is now common among people who build coding agents \citep{anthropic2025harnesses}. None of the names says what the thing contains, where it ends and the model begins, or how two instances of it differ. As a result, agents are compared by the model inside them, and results on one benchmark with one model vary with the system around it (\S\ref{sec:uses}). When such an agent fails, there is no agreed list of places to look.

\paragraph{Position.} We argue three claims.
\begin{enumerate}
  \item \textbf{An LLM agent is a model wrapped in a harness}: $\textit{Agent} = \textit{Harness}(\textit{Model})$. The model maps tokens to a next token. The harness is everything else on the path of a request.
  \item \textbf{Every harness is built from five mechanistic primitives}: \prim{Input}, \prim{Context}, \prim{Model interface}, \prim{Output} and \prim{Control flow}. Each is defined by what it does and by what it never does, so that each piece of a harness has one home.
  \item \textbf{Everything else said about a harness is at one of two other levels}: a \emph{design decision} within a primitive (one tool or forty; a loop or a graph; a repository map or a summarizer), or a \emph{production concern} asked of every primitive (persistence, safety and security, observability, resilience, performance, evaluation). Neither level adds a primitive.
\end{enumerate}

\paragraph{What the vocabulary is for.} It separates what a harness \emph{is} from what a given harness \emph{chose} and from what makes it fit to \emph{run}. With that separation, a failure can be located: was the model shown what it needed (\prim{Context}), did what it asked for happen (\prim{Output}), did the loop stop early or run long (\prim{Control flow}), was the response read correctly (\prim{Model interface}), did the right thing come in (\prim{Input})? Only when all five are ruled out does the failure belong to the model. Two harnesses can be compared decision by decision within the same five primitives, a benchmark result can report the harness that produced it, and a harness can be moved to a new model by changing one primitive.

\paragraph{Evidence.} The definitions come from building a harness from nothing, one primitive at a time, around a fixed model and a fixed task, and observing what each version could not do (Appendix~\ref{app:minimal}). The claim that the five cover existing harnesses is tested by deconstructing six of them from their primary sources at pinned versions: Claude Code, Codex CLI, opencode, pi, Hermes Agent, and a LangGraph workflow (\S\ref{sec:deconstruction}). For each we assign every component on the request path to one primitive, or to one concern together with the primitive it changes, and record every case that was hard to place. The evidence for each assignment is given in Appendix~\ref{app:evidence} and published with the paper.

\paragraph{What this paper is not.} It is a conceptual contribution. It reports no experiment in which a primitive is varied and performance measured, and it does not prove that no sixth primitive could ever be needed. It states what would refute the claim (\S\ref{sec:limitations}) and gives a deconstruction a reader can check.

\paragraph{Contributions.}
\begin{itemize}
  \item A definition of the harness and of its boundary with the model (\S\ref{sec:boundary}).
  \item Five primitives, each with a mechanism, a negative definition, and its design space (\S\ref{sec:primitives}).
  \item Three levels of description, primitive, design decision and production concern, with rules for the cases at their boundaries (\S\ref{sec:levels}).
  \item A deconstruction of six harnesses, with method, selection criteria and published evidence (\S\ref{sec:deconstruction}, Appendix~\ref{app:evidence}).
  \item Uses of the vocabulary: a diagnostic procedure, a reporting schema, and porting (\S\ref{sec:uses}).
  \item The minimal harnesses, one primitive at a time (Appendix~\ref{app:minimal}).
\end{itemize}
